\hypertarget{testing}{%
\section{Testing}\label{testing}}

Testing here is kept separate from evaluation. This section checks that
the components behave as designed. Section~\ref{evaluation-and-discussion} is where the assembled system
is compared against baselines instead. Two complementary strategies were
used throughout development. Both were chosen to fit a project built by
one developer on a fifteen-week timeline.

Two further things shaped how this section and Section~\ref{evaluation-and-discussion} are
structured, not just what they test.

First, this project follows Krishnamachari's (2026) checklist for
defensible empirical results. Several items on that checklist map
directly onto Section~\ref{evaluation-and-discussion}'s design. Strong, named baselines are used
instead of a single straw-man comparison. Five core baselines plus
eight literature re-implementations are run against the proposed
system (Section~\ref{evaluation-strategy}). Multiple regimes are tested instead of one
operating point, using three traffic scenarios. And the report
states both success and failures of a method. Section~\ref{limitations} discloses
the degraded-network over-escalation limitation rather than smoothing
it over.

Second, self-check correctness and live-traffic testing are kept as
genuinely separate activities in this section. A passing self-check
suite shows that a component behaves as designed in isolation, but
whether the deployed system behaves the same way under real network
calls, real models, and real timing is a separate question that
self-checks alone cannot answer. Section~\ref{live-integration-tests}
covers what answering that second question required, and what it found
that the self-check suite alone had not.

\hypertarget{self-check-convention}{%
\subsection{The self-check convention}\label{self-check-convention}}

Every module in \textbf{devmind/} that contains non-trivial logic (a
branch, a loop, a numeric threshold, a state machine) carries its own
\textbf{demo()} function. It can be run directly (\textbf{python -m
devmind.\textless module\textgreater}), and it asserts the behaviours
that would break silently if the logic regressed. For example,
\textbf{edge.py}'s self-check verifies that
\textbf{compute\_calibration\_delta()} returns a signed correction.
\textbf{agent.py}'s self-check verifies that
\textbf{is\_ood()} needs two anomalous slots, not one, and that a stale
\textbf{last\_fallback\_reason} is cleared on the next non-fallback
decision. \textbf{orchestrator.py}'s self-check verifies, among other
things, that a client with an explicit threshold override is not moved
by \textbf{recalibrate\_thresholds()}, and that the LLM governance
review can escalate a decision but never downgrade one.

There is no pytest suite with fixtures and mocks. Self-checks are
introduced as a small, readable, standalone script. Each of the self-checks
is like a small readable standalone script. The test uses the actual methods
and real inputs, even very minimal ones at times, and when it fails, it produces
a clear and meaningful error that is not just stating the failure itself but also
pointing out the reasoning why the invariant matters. Every module touched during
this project's latervsprints carries one of these self-checks (\textbf{edge.py},
\textbf{dataset.py}, \textbf{cascade.py}, \textbf{agent.py},
\textbf{orchestrator.py}, \textbf{diagnosis.py},
\textbf{orchestrator\_trainer.py},
\textbf{gateway/orchestrator\_app.py}), and none of them makes a real
network call. External dependencies, HuggingFace Hub and Ollama, are
stood in for with small fake objects that match the real interface. For
example, \textbf{orchestrator\_trainer.py}'s \textbf{\_FakeModel}
matches \textbf{DistilBERTEdge}/\textbf{BERTLargeCloud}'s
\textbf{.predict(text, true\_label)} signature without calling a real
model. So the whole suite runs in well under a minute, and it can be
re-run after any change without needing a live model or a network
connection.

Two bugs were found and fixed directly through this discipline.
\textbf{compute\_calibration\_delta()} in \textbf{edge.py}
was subtracting an unsigned gap from raw confidence,
instead of applying the signed, temperature-scaled correction. And
\textbf{TextClassificationDataset.shuffle\_train()} in
\textbf{dataset.py} was shuffling a throwaway filtered copy of the
training splits instead of the samples actually used for training. Both
bugs were caught while writing the self-check for each module, during a
full read-through of the codebase. Neither was caught by a failing test
that had already shipped.

\hypertarget{live-integration-tests}{%
\subsection{Live integration tests}\label{live-integration-tests}}

Self-checks deliberately avoid real models, HTTP connections, and LLM
calls. So a second layer of testing exercised the system with all three
present, for the parts where a fake cannot stand in for the real
failure surface.

\begin{itemize}
\item \textbf{Multi-region live traffic test} (documented in
\textbf{evaluation/multi-region-live-traffic-test-2026-08-05.md}). The
three Docker services (edge gateway, cloud pod, orchestrator dashboard)
were deployed to real GCE VMs across three regions via
\textbf{deploy/gcp-up.sh}, and driven with real traffic. This is what
surfaced a genuine gap that no unit-level test had caught. The cloud
dispatch path in \textbf{cascade.py} had no \textbf{try/except} around
the network call, unlike the edge path, which does. This has since been
fixed. \textbf{\_dispatch()} now retries the cloud call once inside a
\textbf{try/except}, then falls back to the already-computed edge
result if it still fails (Section~\ref{design-development-and-implementation}, Cascade Controller). That is
exactly the class of gap a live multi-region test can surface, and a
self-check running against local fakes cannot.
\item \textbf{LLM diagnosis capability}. This was tested end-to-end against a
real local Ollama instance (\textbf{gemma4:12b-it-qat}), not just the
scripted stub used in self-checks. A synthetic \textbf{client\_babcock}
distress scenario produced a coherent, correctly-parsed diagnosis in
roughly 2.5 seconds once the model was warm. A first, cold-Ollama
attempt timed out at the then-configured 15 seconds. That is why the
provider's default timeout was raised to 20 seconds.
\item \textbf{Websocket dashboard}. This was tested with a real
\textbf{uvicorn} server, a real websocket client, and temporary
action-log files, not just the in-memory fakes (\textbf{FakeWS},
\textbf{DeadWS}, \textbf{FakeMonitor}) used in
\textbf{orchestrator\_app.py}'s self-check.
\item \textbf{LLM governance review}. This was live-tested against real
Ollama, using a real client's numbers (Section~\ref{design-development-and-implementation}'s ``LLM governance
review''). It correctly escalated a \textbf{fine\_tune} decision to
\textbf{train\_new}, given a genuine large SLA-violation gap, and gave
a coherent justification for doing so.
\item \textbf{Multi-region Kubernetes deployment} (documented in
\textbf{evaluation/multi-region-gke-live-traffic-test-2026-09-02.md}).
The three services were deployed via real \textbf{kubectl}-managed
\textbf{Deployment}/\textbf{Service} manifests, onto one small GKE
cluster per region. A single cluster cannot span the $\approx$250ms RTT to the
Sydney region, so cross-region calls stay at the HTTP layer, the same
as in the Docker path. This deployment ran the classifier-retrained
policy (Section~\ref{misc-classifier-refit}), and handled 9,835 requests with zero errors and
zero pod restarts across the full session.

The live escalation rate tracks the matching offline condition closely.
For \textbf{client\_babcock} it was 2.5\% live against 3.5\% offline,
and the other clients were within 0.2--0.7 points. These differences
are explained by real network variance, not a regression. The matching
offline condition uses the retrained policy paired with the untrained
heuristic, not the trained classifier. \textbf{gateway/app.py} only
reads a fitted classifier when \textbf{DEVMIND\_MISC\_CLASSIFIER\_PATH}
is set, and this deployment deliberately left it unset. That classifier
was fitted on the simulation's synthetic stress distribution, and it
has never been validated against real \textbf{psutil} readings (Section~\ref{misc-classifier-refit}). So the live gateway and the offline simulation default to
different classifiers. This is a disclosed scoping decision, not an
oversight, and Section~\ref{misc-classifier-refit}'s comparison is what makes it a defensible
one. The retrained policy turns out to be nearly insensitive to which
of the two feeds it.

This deployment also ran Prometheus and Grafana as real Kubernetes
workloads, scraping \textbf{prometheus\_client} metrics across all
three regions. It surfaced two further gaps that a self-check running
against local fakes cannot catch. First,
\textbf{prometheus\_client}'s default \textbf{Histogram} buckets are
calibrated for second-scale durations, and they silently degraded
millisecond-scale latency quantiles until explicit buckets were set.
Second, the orchestrator's live-monitoring view tails a local log file,
and it needed an explicit HTTP-forwarding path
(\textbf{CascadeController}'s \textbf{action\_log\_url}) once the
gateway and orchestrator became separate pods rather than processes
sharing one host. Both problems were fixed and verified against live
cross-region traffic.

One constraint remains, and it is disclosed rather than hidden. The GCP
project's global \textbf{IN\_USE\_ADDRESSES} quota is 8, and this
deployment's own 8 services already hold all of it. So Jaeger cannot
also acquire an external IP. The pod runs, but its UI and OTLP endpoint
are not externally reachable under this project's current quota.
\end{itemize}

\hypertarget{what-was-not-tested}{%
\subsection{What was not tested}\label{what-was-not-tested}}

The single-node EC2+Minikube deployment path
(\textbf{deploy/ec2-minikube-setup.sh}) is checked into the repository.
It looks plausible on inspection, but it has no live-test report to
back it up. The 3-region GKE deployment above is the Kubernetes path
that does have live-test evidence, using genuine multi-cluster
Kubernetes rather than a single local cluster.

The learned orchestrator meta-policy (Section~\ref{design-development-and-implementation}) was tested in one
limited sense. Its integration code, safety fallback, and training loop
all pass self-checks, and were run against real training data. But it
was never tested in the sense of producing a trusted live decision,
because it never converged. That is reported here as a failure, not as
a gap in test coverage.